\begin{tikzfig}
\begin{adjustbox}{max width=\linewidth}\begin{tikzpicture}[fig, x=1.05cm, y=1cm]
% дорожки
\node[font=\sffamily\small\bfseries, anchor=east, align=right] at (-0.15,2.2) {Ваша \\ программа};
\node[font=\sffamily\small\bfseries, anchor=east] at (-0.15,0.45) {Автопилот};
% программа
\fill[Orange!10, draw=Accent, line width=0.8pt, rounded corners=3pt] (0,1.7) rectangle (2,2.7);
\fill[Orange!10, draw=Accent, line width=0.8pt, rounded corners=3pt] (2.05,1.7) rectangle (5,2.7);
\fill[Orange!10, draw=Accent, line width=0.8pt, rounded corners=3pt] (5.05,1.7) rectangle (5.55,2.7);
\fill[Success!10, draw=Success, line width=0.8pt, rounded corners=3pt] (5.6,1.7) rectangle (9,2.7);
\fill[Orange!10, draw=Accent, line width=0.8pt, rounded corners=3pt] (9.05,1.7) rectangle (12,2.7);
\node[font=\ttfamily\small] at (1,2.2) {arm()};
\node[font=\ttfamily\small] at (3.5,2.2) {takeoff()};
\node[font=\sffamily\footnotesize, text width=3.3cm, align=center] at (7.3,2.2) {код работает дальше, например читает датчики};
\node[font=\ttfamily\small] at (10.5,2.2) {land()};
% автопилот
\foreach \a/\b/\t in {0/2/моторы, 2.05/5/взлёт, 5.05/9/перелёт к точке, 9.05/12/посадка} {
  \fill[Paper, draw=Muted, line width=0.8pt, rounded corners=3pt] (\a,0) rectangle (\b,0.9);
  \node[font=\sffamily\small] at ({(\a+\b)/2},0.45) {\t};
}
% пояснение
\node[font=\sffamily\footnotesize, anchor=center] at (3.5,1.3) {блокирует: ждёт события};
\draw[fcmd] (5.3,3.4) -- (5.3,2.75);
\node[font=\sffamily\small, anchor=south] at (5.3,3.4) {\texttt{go\_to\_local\_point()} возвращается сразу};
\draw[faxis] (0,-0.5) -- (12.4,-0.5) node[right, font=\sffamily\small] {время};
\end{tikzpicture}\end{adjustbox}
\captionof{figure}{Блокирующие команды (\texttt{arm}, \texttt{takeoff}, \texttt{land}) ждут события автопилота. Полётные команды (\texttt{go\_to\_*}, \texttt{set\_manual\_speed*}, \texttt{rtl}) только отправляют команду, и программа продолжает работу, пока дрон летит.}\label{fig:sdk2-blocking}
\end{tikzfig}
